\section*{ПРИЛОЖЕНИЕ 1. ФРАГМЕНТЫ ИСХОДНОГО КОДА} \label{application_1}
\addcontentsline{toc}{section}{ПРИЛОЖЕНИЕ 1. Фрагменты исходного кода}
\pagestyle{empty}

\setcounter{figure}{0}

В настоящем приложении приведены ключевые фрагменты исходного
кода реализованного сервиса, соответствующие моделям, описанным
в главе~2 диссертации. Полный исходный код опубликован в репозиториях
\texttt{eduquiz-backend} (серверная часть на языке Go) и
\texttt{eduquiz-flutter} (кроссплатформенный клиент на Flutter).

\subsection*{П.1.1. Реализация алгоритма SM\textendash 2}

Реализация чистая, не зависящая от инфраструктуры; используется
как обработчиками REST-запросов, так и в юнит-тестах. Файл
\texttt{internal/services/sm2.go} серверной части.

\begin{lstlisting}[caption={Алгоритм интервального повторения SM-2},label=lst:sm2-go]
// Sm2State хранит параметры повторения карточки для пары (студент, вопрос).
type Sm2State struct {
    EF             float64   // easiness factor; init=2.5, min=1.3
    IntervalDays   int       // текущий интервал повторения в днях
    Repetitions    int       // число последовательных успешных повторений
    DueDate        time.Time // дата следующего планового повторения
    LastReviewedAt time.Time // момент последней оценки
}

// Sm2Update применяет ответ с оценкой качества q in [0,5].
// Возвращает обновлённое состояние; не мутирует s.
func Sm2Update(s Sm2State, q int, now time.Time) Sm2State {
    if q < 0 { q = 0 }
    if q > 5 { q = 5 }

    out := s
    if q < 3 {
        out.Repetitions = 0
        out.IntervalDays = 1
    } else {
        switch s.Repetitions {
        case 0: out.IntervalDays = 1
        case 1: out.IntervalDays = 6
        default: out.IntervalDays = ceilInt(float64(s.IntervalDays) * s.EF)
        }
        out.Repetitions = s.Repetitions + 1
    }

    delta := 0.1 - float64(5-q)*(0.08+float64(5-q)*0.02)
    out.EF = s.EF + delta
    if out.EF < 1.3 {
        out.EF = 1.3
    }

    out.LastReviewedAt = now
    out.DueDate = now.Add(time.Duration(out.IntervalDays) * 24 * time.Hour)
    return out
}
\end{lstlisting}

\subsection*{П.1.2. Автоматическое формирование оценки качества ответа}

Оценка $q$ для алгоритма SM\textendash 2 формируется автоматически
по корректности и времени ответа без явного запроса у обучающегося.
Фрагмент из обработчика ответа в файле \texttt{internal/handlers/rooms.go}.

\begin{lstlisting}[caption={Формирование q-score из (корректность $\times$ время)},label=lst:qscore-go]
// sm2QScore выводит q in [0,5] из факта правильности и скорости ответа.
//   correct + <=½T -> 5;  correct + <=T -> 4;  correct + >T -> 3;
//   incorrect      -> 2;  таймаут        -> 1.
func sm2QScore(correct bool, elapsedMs, timeLimitMs int) int {
    if !correct {
        if timeLimitMs > 0 && elapsedMs >= timeLimitMs {
            return 1
        }
        return 2
    }
    if timeLimitMs <= 0 {
        return 4
    }
    if elapsedMs*2 <= timeLimitMs {
        return 5
    }
    if elapsedMs <= timeLimitMs {
        return 4
    }
    return 3
}
\end{lstlisting}

\subsection*{П.1.3. Расчёт очков опыта (XP)}

Реализация формулы~\eqref{eq:xp} из раздела~2.4. Файл
\texttt{internal/services/xp.go} серверной части.

\begin{lstlisting}[caption={Расчёт XP за один ответ},label=lst:xp-go]
type XPParams struct {
    Alpha     float64 // штраф за медлительность; 0.5 по умолчанию
    MinFactor float64 // минимум скоростного множителя; 0.0 по умолчанию
}

// XPAward вычисляет XP за один ответ.
//   correct      — правильный ли ответ;
//   difficulty   — сложность вопроса 1..5 (вне диапазона клампится);
//   elapsedMs    — время ответа в миллисекундах;
//   timeLimitMs  — лимит времени на вопрос (0 = без лимита).
func XPAward(correct bool, difficulty int, elapsedMs, timeLimitMs int, p XPParams) int {
    if !correct {
        return 0
    }
    if difficulty < 1 { difficulty = 1 }
    if difficulty > 5 { difficulty = 5 }

    speed := 1.0
    if timeLimitMs > 0 {
        ratio := float64(elapsedMs) / float64(timeLimitMs)
        if ratio > 1 {
            ratio = 1
        }
        speed = 1.0 - p.Alpha*ratio
        if speed < p.MinFactor {
            speed = p.MinFactor
        }
    }

    xp := float64(difficulty) * speed * 10.0
    return int(xp + 0.5)
}

// LevelFromXP возвращает уровень по правилу L = floor(sqrt(XP/100)).
func LevelFromXP(xp int) int {
    const k = 100
    if xp <= 0 {
        return 0
    }
    l := 0
    for {
        next := (l + 1) * (l + 1) * k
        if next > xp {
            return l
        }
        l++
    }
}
\end{lstlisting}

\subsection*{П.1.4. Расчёт коэффициента $\alpha$ Кронбаха и статистик пунктов}

Реализация формулы~\eqref{eq:cronbach}, индекса трудности $p_i$
и точечно-бисериальной корреляции $r_{pb,i}$. Файл
\texttt{internal/services/analytics.go}.

\begin{lstlisting}[caption={Расчёт $\alpha$ Кронбаха и статистик пунктов},label=lst:analytics-go]
// AnalyzeBank вычисляет alpha + статистики каждого пункта по матрице
// ответов responses[participant][item] in {0,1}.
func AnalyzeBank(responses [][]int) BankAnalytics {
    n := len(responses)
    if n < 2 { return BankAnalytics{Respondents: n} }
    k := len(responses[0])
    if k < 2 { return BankAnalytics{Respondents: n, ItemCount: k} }

    totals   := make([]float64, n)
    itemSums := make([]float64, k)
    for i := 0; i < n; i++ {
        for j := 0; j < k; j++ {
            x := float64(responses[i][j])
            totals[i]   += x
            itemSums[j] += x
        }
    }
    totalMean := mean(totals)
    totalVar  := variancePopulation(totals, totalMean)

    // Дисперсия каждого пункта; для бинарных = p(1-p).
    sumItemVar := 0.0
    for j := 0; j < k; j++ {
        p := itemSums[j] / float64(n)
        sumItemVar += p * (1 - p)
    }

    // alpha Кронбаха: (k/(k-1)) * (1 - sum sigma^2_i / sigma^2_total).
    var alpha float64
    if totalVar > 0 {
        alpha = float64(k) / float64(k-1) * (1 - sumItemVar/totalVar)
    }

    // Per-item: p-value и r_pb (точечно-бисериальная корреляция).
    totalStd := math.Sqrt(totalVar)
    items := make([]ItemStats, k)
    for j := 0; j < k; j++ {
        p := itemSums[j] / float64(n)
        var sumRight, sumWrong float64
        var nRight, nWrong int
        for i := 0; i < n; i++ {
            if responses[i][j] == 1 {
                sumRight += totals[i]; nRight++
            } else {
                sumWrong += totals[i]; nWrong++
            }
        }
        var rpb float64
        if nRight > 0 && nWrong > 0 && totalStd > 0 {
            mRight := sumRight / float64(nRight)
            mWrong := sumWrong / float64(nWrong)
            rpb = (mRight - mWrong) / totalStd * math.Sqrt(p*(1-p))
        }
        items[j] = ItemStats{PValue: p, Discrimination: rpb}
    }
    return BankAnalytics{Respondents: n, ItemCount: k,
        CronbachAlpha: alpha, Items: items}
}
\end{lstlisting}

\subsection*{П.1.5. Цветовая интерпретация значения $\alpha$ на клиенте}

Фрагмент из файла \texttt{lib/features/teacher/bank\_analytics\_screen.dart}
клиентской части (Flutter). Цветовая интерпретация значения $\alpha$
и флажки качества пунктов используются преподавателем для
итеративного улучшения банка вопросов.

\begin{lstlisting}[caption={Интерпретация $\alpha$ Кронбаха в клиенте},label=lst:alpha-dart]
String alphaInterpretation(double a) {
  if (a >= 0.9) return 'Отлично — но проверьте дублирование пунктов';
  if (a >= 0.8) return 'Хорошо — банк надёжен';
  if (a >= 0.7) return 'Приемлемо — удалите слабые пункты';
  if (a >= 0.6) return 'Сомнительно — требуется доработка';
  return 'Низко — банк требует серьёзной переработки';
}

Color alphaColor(double a, ColorScheme scheme) {
  if (a >= 0.8) return scheme.success;
  if (a >= 0.7) return scheme.primary;
  if (a >= 0.6) return scheme.warning;
  return scheme.danger;
}
\end{lstlisting}

\subsection*{П.1.6. Бутстреп-оценка доверительного интервала
размера эффекта Коэна}

Фрагмент пайплайна вычислительного контура обработки результатов
(подраздел~3.9 главы~3), реализующий бутстреп-оценку доверительного
интервала размера эффекта Коэна $d$ методом BCa (bias-corrected
and accelerated). Файл \texttt{experiment/lib/effect\_size.py}.

\begin{lstlisting}[language=Python,caption={Бутстреп-CI для $d$ Коэна (BCa)},label=lst:bootstrap-d-py]
import numpy as np
from scipy.stats import norm

def cohen_d(x: np.ndarray, y: np.ndarray) -> float:
    """Standardised mean difference using pooled SD."""
    nx, ny = len(x), len(y)
    s_p = np.sqrt(((nx - 1) * x.var(ddof=1) +
                   (ny - 1) * y.var(ddof=1)) / (nx + ny - 2))
    return (x.mean() - y.mean()) / s_p

def bootstrap_bca_ci(
    x_exp: np.ndarray, x_ctrl: np.ndarray,
    B: int = 10_000, alpha: float = 0.05, rng_seed: int = 42,
) -> tuple[float, float, float]:
    rng = np.random.default_rng(rng_seed)
    d_hat = cohen_d(x_exp, x_ctrl)

    # 1) Bootstrap distribution.
    d_b = np.empty(B)
    for b in range(B):
        xe = rng.choice(x_exp, size=len(x_exp), replace=True)
        xc = rng.choice(x_ctrl, size=len(x_ctrl), replace=True)
        d_b[b] = cohen_d(xe, xc)

    # 2) Bias correction.
    z0 = norm.ppf(np.mean(d_b < d_hat))

    # 3) Acceleration via jackknife.
    pooled = np.concatenate([x_exp, x_ctrl])
    jk = np.empty(len(pooled))
    for i in range(len(pooled)):
        mask = np.ones(len(pooled), dtype=bool); mask[i] = False
        xe = pooled[:len(x_exp)] if i >= len(x_exp) else np.delete(x_exp, i)
        xc = pooled[len(x_exp):] if i < len(x_exp) else np.delete(x_ctrl, i - len(x_exp))
        jk[i] = cohen_d(xe, xc)
    jk_mean = jk.mean()
    num = np.sum((jk_mean - jk) ** 3)
    den = 6.0 * (np.sum((jk_mean - jk) ** 2) ** 1.5)
    a = num / den if den > 0 else 0.0

    # 4) BCa quantiles.
    z_lo, z_hi = norm.ppf(alpha / 2), norm.ppf(1 - alpha / 2)
    a_lo = norm.cdf(z0 + (z0 + z_lo) / (1 - a * (z0 + z_lo)))
    a_hi = norm.cdf(z0 + (z0 + z_hi) / (1 - a * (z0 + z_hi)))
    lo, hi = np.quantile(d_b, [a_lo, a_hi])
    return d_hat, lo, hi
\end{lstlisting}

\subsection*{П.1.7. Бутстреп-оценка доверительного интервала
коэффициента \texorpdfstring{$\alpha$}{alpha} Кронбаха}

Реализация бутстрепа по строкам матрицы ответов $U \in \{0,1\}^{n \times k}$.
Файл \texttt{experiment/lib/reliability.py}.

\begin{lstlisting}[language=Python,caption={Бутстреп-CI для $\alpha$ Кронбаха},label=lst:bootstrap-alpha-py]
import numpy as np

def cronbach_alpha(U: np.ndarray) -> float:
    """Cronbach alpha for binary item-response matrix U (n x k)."""
    n, k = U.shape
    item_var = U.var(axis=0, ddof=1).sum()
    total_var = U.sum(axis=1).var(ddof=1)
    if total_var <= 0:
        return float("nan")
    return (k / (k - 1)) * (1 - item_var / total_var)

def alpha_bootstrap_ci(
    U: np.ndarray, B: int = 10_000, alpha_level: float = 0.05,
    rng_seed: int = 42,
) -> tuple[float, float, float]:
    rng = np.random.default_rng(rng_seed)
    n = U.shape[0]
    alpha_hat = cronbach_alpha(U)
    alpha_b = np.empty(B)
    for b in range(B):
        idx = rng.integers(0, n, size=n)  # resample rows with replacement
        alpha_b[b] = cronbach_alpha(U[idx])
    lo, hi = np.quantile(alpha_b, [alpha_level / 2, 1 - alpha_level / 2])
    return alpha_hat, lo, hi
\end{lstlisting}

\subsection*{П.1.8. Калибровка параметров банка по
двухпараметрической модели IRT}

Скелет MML-калибровки $(a_j, b_j)$ через интерфейс
библиотеки \texttt{py-irt} (или эквивалентного моста к R-пакету
\texttt{mirt}); ниже приведён фрагмент c упрощённым EM-итерированием
для иллюстрации принципиальной структуры алгоритма. Полноценная
устойчивая реализация использует библиотеку \texttt{py-irt}.

\begin{lstlisting}[language=Python,caption={IRT-калибровка $(a_j, b_j)$ (схема)},label=lst:irt-calib-py]
import numpy as np
from scipy.special import expit
from scipy.optimize import minimize

def irt_2pl_likelihood(params: np.ndarray, U: np.ndarray,
                       theta_grid: np.ndarray, prior: np.ndarray) -> float:
    """Marginal log-likelihood for 2PL IRT over a quadrature grid."""
    n, k = U.shape
    a = params[:k]
    b = params[k:]
    # P_{ij}(theta) per quadrature node
    Theta = theta_grid[:, None, None]                  # (G, 1, 1)
    P = expit(a[None, None, :] * (Theta - b[None, None, :]))  # (G, 1, k)
    # Per-respondent likelihood over quadrature
    LL_i_theta = (U[None, :, :] * np.log(P + 1e-12)
                  + (1 - U[None, :, :]) * np.log(1 - P + 1e-12)).sum(axis=2)
    marginal = (np.exp(LL_i_theta) * prior[:, None]).sum(axis=0)
    return -np.log(marginal + 1e-12).sum()

def calibrate_2pl(U: np.ndarray, n_grid: int = 41) -> dict:
    n, k = U.shape
    theta_grid = np.linspace(-4, 4, n_grid)
    prior = np.exp(-theta_grid ** 2 / 2) / np.sqrt(2 * np.pi)
    prior /= prior.sum()

    x0 = np.concatenate([np.ones(k), np.zeros(k)])  # a_j=1, b_j=0
    res = minimize(
        irt_2pl_likelihood, x0, args=(U, theta_grid, prior),
        method="L-BFGS-B",
        bounds=[(0.3, 4.0)] * k + [(-4.0, 4.0)] * k,
    )
    a, b = res.x[:k], res.x[k:]
    return {"a": a, "b": b, "converged": res.success}
\end{lstlisting}

\subsection*{П.1.9. Обучение валидити-классификатора}

Фрагмент обучения бинарного классификатора валидности атрибуции
ответа на размеченной выборке. Файл
\texttt{experiment/lib/validity\_classifier.py}.

\begin{lstlisting}[language=Python,caption={Валидити-классификатор: обучение и калибровка},label=lst:validity-py]
import numpy as np
import pandas as pd
from sklearn.linear_model import LogisticRegression
from sklearn.ensemble import GradientBoostingClassifier
from sklearn.calibration import CalibratedClassifierCV
from sklearn.model_selection import StratifiedKFold, cross_val_score

FEATURES = [
    "elapsed_z_score",      # z-score времени ответа vs. профиль
    "blur_ratio",           # доля внефокусного времени
    "paste_ratio",          # доля вставленных символов
    "iki_dev",              # отклонение биометрии набора
    "tab_switches",         # частота переключений вкладок
    "consistency_score",    # консистентность с уточняющими вопросами
]

def train_validity_clf(df: pd.DataFrame, model: str = "logreg"):
    """Train binary classifier: y = 1 if ответ был самостоятельным."""
    X = df[FEATURES].to_numpy()
    y = df["is_independent"].to_numpy().astype(int)

    if model == "logreg":
        base = LogisticRegression(class_weight="balanced", max_iter=200)
    elif model == "gbm":
        base = GradientBoostingClassifier(
            n_estimators=200, max_depth=3, learning_rate=0.05,
        )
    else:
        raise ValueError(f"unknown model {model!r}")

    # Calibrated probability output (Platt scaling on held-out fold).
    clf = CalibratedClassifierCV(base, cv=5, method="sigmoid")

    cv_auc = cross_val_score(
        clf, X, y, cv=StratifiedKFold(5, shuffle=True, random_state=42),
        scoring="roc_auc",
    )
    clf.fit(X, y)
    return clf, {"cv_auc_mean": cv_auc.mean(), "cv_auc_std": cv_auc.std()}

def integral_score(theta_hat: float, validity_prob: float) -> float:
    """Final integrated knowledge score: Score = theta_hat * validity."""
    return theta_hat * float(np.clip(validity_prob, 0.0, 1.0))
\end{lstlisting}

\newpage
\section*{ПРИЛОЖЕНИЕ 2. СКРИНШОТЫ ИНТЕРФЕЙСА} \label{application_2}
\addcontentsline{toc}{section}{ПРИЛОЖЕНИЕ 2. Скриншоты интерфейса}

Скриншоты основных экранов клиентского приложения, иллюстрирующие
функциональность, описанную в главе~2 диссертации.

\subsection*{П.2.1. Экран входа и регистрации}

\image{ui-auth.png}{Экран входа и регистрации (Web)}{fig:ui-auth}{0.7}

\subsection*{П.2.2. Кабинет преподавателя}

Список курсов, экран банка вопросов, психометрическая карточка
($\alpha$ Кронбаха и статистики пунктов).

\image{ui-teacher-courses.png}{Кабинет преподавателя: список курсов}{fig:ui-teacher-courses}{0.7}
\image{ui-bank-editor.png}{Редактор банка вопросов}{fig:ui-bank-editor}{0.7}
\image{ui-bank-analytics.png}{Психометрическая карточка банка вопросов}{fig:ui-bank-analytics}{0.7}

\subsection*{П.2.3. Кабинет студента}

Ввод кода сессии, экран ответа на вопрос, итоговый экран сессии.

\image{ui-student-home.png}{Кабинет студента: ввод кода сессии}{fig:ui-student-home}{0.7}
\image{ui-room-active.png}{Активный вопрос интерактивной сессии}{fig:ui-room-active}{0.7}
\image{ui-room-finished.png}{Итоговый экран сессии}{fig:ui-room-finished}{0.7}

\subsection*{П.2.4. Экран интервального повторения SM\textendash 2}

\image{ui-sm2.png}{Экран индивидуального повторения по SM-2}{fig:ui-sm2}{0.7}

\subsection*{П.2.5. Профиль и игровая прогрессия}

XP, уровень, серия дней, гистограмма XP за последние 14 дней,
каталог бейджей.

\image{ui-profile.png}{Профиль студента: XP, уровень, бейджи}{fig:ui-profile}{0.7}

\subsection*{П.2.6. Лидерборд интерактивной сессии}

\image{ui-leaderboard.png}{Лидерборд live-сессии}{fig:ui-leaderboard}{0.7}

\newpage
\section*{ПРИЛОЖЕНИЕ 3. ПРОТОКОЛЫ ПЕДАГОГИЧЕСКОГО ЭКСПЕРИМЕНТА} \label{application_3}
\addcontentsline{toc}{section}{ПРИЛОЖЕНИЕ 3. Протоколы педагогического эксперимента}

В настоящем приложении приведены формы протоколов педагогического
эксперимента в соответствии с методикой, описанной в разделах~3.5
и~3.6 диссертации. Заполненные протоколы прилагаются к работе
после завершения апробации в текущем семестре.

\subsection*{П.3.1. Форма анкеты участника эксперимента}

Анкета заполняется студентом перед началом эксперимента
с обезличенным идентификатором (присваивается экспериментатором).

\begin{itemize}
    \item Идентификатор участника: \dotfill
    \item Группа: \,контрольная\,/\,экспериментальная (нужное обвести)
    \item Курс обучения: \dotfill
    \item Опыт использования интерактивных образовательных приложений
          (Kahoot, Quizlet, Mentimeter и др.):
          \,нет\,/\,эпизодически\,/\,регулярно\,/\,постоянно
    \item Согласие на обработку обезличенных результатов в составе
          выборки эксперимента: \,да\,/\,нет
\end{itemize}

\subsection*{П.3.2. Структура pre-/post-/retention-теста}

Тестовый инструмент: 18~пунктов с одиночным выбором по содержанию
курса \textit{«Надёжность»} (терминология теории надёжности,
показатели и формулы, схемы резервирования). Каждый тест
проводится в одинаковых условиях и в течение одинакового времени.
\begin{itemize}
    \item Pre-test --- начало эксперимента (до вмешательства);
    \item Post-test --- завершение цикла обучающих сессий;
    \item Retention-test --- через 7--14 суток после post-test;
\end{itemize}

\subsection*{П.3.3. Форма обезличенной таблицы результатов}

Таблица заполняется по итогам каждого замера и используется как вход
для программного пайплайна статистического анализа (раздел~3.7).
Формат данных соответствует файлу
\texttt{experiment/data/scores\_sample.csv}: \texttt{participant\_id},
\texttt{group} (\texttt{control} / \texttt{experimental}),
\texttt{pre\_score}, \texttt{post\_score}, \texttt{retention\_score},
\texttt{sessions\_count}.

Дополнительно собирается матрица бинарных ответов по пунктам
(\texttt{participant\_id}, \texttt{group}, \texttt{q1}, \dots, \texttt{q18})
для расчёта $\alpha$ Кронбаха и точечно-бисериальной корреляции;
формат соответствует файлу \texttt{experiment/data/items\_sample.csv}.

\newpage
\section*{ПРИЛОЖЕНИЕ 4. СПИСОК СОКРАЩЕНИЙ И ОБОЗНАЧЕНИЙ} \label{application_4}
\addcontentsline{toc}{section}{ПРИЛОЖЕНИЕ 4. Список сокращений и обозначений}

\subsection*{П.4.1. Аббревиатуры}

\begin{description}
    \item[API] --- Application Programming Interface, программный интерфейс приложения.
    \item[CITEXT] --- Case-Insensitive Text, регистронезависимый текстовый тип PostgreSQL.
    \item[CSRF] --- Cross-Site Request Forgery, межсайтовая подделка запроса.
    \item[CSV] --- Comma-Separated Values, формат данных с разделителями-запятыми.
    \item[ENUM] --- Enumeration, перечислимый тип данных.
    \item[FSM] --- Finite-State Machine, конечный автомат.
    \item[GIFT] --- General Import Format Template, текстовый формат банка вопросов Moodle.
    \item[JSON] --- JavaScript Object Notation, формат обмена данными.
    \item[JWT] --- JSON Web Token, защищённый токен в формате JSON.
    \item[LMS] --- Learning Management System, система управления обучением.
    \item[MTBF] --- Mean Time Between Failures, средняя наработка между отказами.
    \item[MTTF] --- Mean Time To Failure, средняя наработка до отказа.
    \item[MTTR] --- Mean Time To Repair, среднее время восстановления.
    \item[OWASP] --- Open Web Application Security Project, открытый проект по безопасности веб-приложений.
    \item[RBD] --- Reliability Block Diagram, структурная схема надёжности.
    \item[REST] --- Representational State Transfer, архитектурный стиль распределённых систем.
    \item[RPO] --- Recovery Point Objective, целевая точка восстановления.
    \item[RTO] --- Recovery Time Objective, целевое время восстановления.
    \item[SLA] --- Service Level Agreement, соглашение об уровне обслуживания.
    \item[SM-2] --- SuperMemo Method, Algorithm 2, алгоритм интервального повторения.
    \item[SQL] --- Structured Query Language, язык структурированных запросов.
    \item[TLS] --- Transport Layer Security, протокол защиты транспортного уровня.
    \item[UML] --- Unified Modeling Language, унифицированный язык моделирования.
    \item[UUID] --- Universally Unique Identifier, универсально-уникальный идентификатор.
    \item[XSS] --- Cross-Site Scripting, межсайтовое выполнение сценариев.
    \item[ВБР] --- вероятность безотказной работы.
    \item[ВКР] --- выпускная квалификационная работа.
    \item[ВУЗ] --- высшее учебное заведение.
    \item[ГОСТ] --- государственный стандарт.
    \item[ИБ] --- информационная безопасность.
    \item[КГ] --- контрольная группа педагогического эксперимента.
    \item[НИР] --- научно-исследовательская работа.
    \item[НИУ] --- национальный исследовательский университет.
    \item[ОС] --- операционная система.
    \item[ПО] --- программное обеспечение.
    \item[РФ] --- Российская Федерация.
    \item[СУБД] --- система управления базами данных.
    \item[ЭГ] --- экспериментальная группа педагогического эксперимента.
\end{description}

\subsection*{П.4.2. Основные обозначения}

\begin{description}
    \item[$P(t)$] --- вероятность безотказной работы за время~$t$.
    \item[$\lambda$] --- интенсивность отказов.
    \item[$T_\text{в}$] --- среднее время восстановления узла.
    \item[$k_\text{г}$] --- коэффициент готовности.
    \item[$\alpha$] --- в формуле XP --- коэффициент штрафа за медлительность; в статистическом анализе --- уровень значимости и коэффициент $\alpha$ Кронбаха.
    \item[$EF$] --- easiness factor, коэффициент лёгкости в алгоритме SM\textendash 2.
    \item[$I$] --- интервал повторения в днях (SM\textendash 2).
    \item[$q$] --- оценка качества ответа в алгоритме SM\textendash 2, $q\in[0,5]$.
    \item[$d$] --- оценка сложности вопроса, $d\in[1,5]$.
    \item[$XP$] --- очки опыта, накапливаемые обучающимся.
    \item[$L$] --- уровень обучающегося в системе игровой прогрессии.
    \item[$p_i$] --- индекс трудности $i$-го пункта тестового инструмента.
    \item[$r_{pb,i}$] --- точечно-бисериальная корреляция $i$-го пункта с суммарным баллом.
    \item[$d_\text{Cohen}$] --- размер эффекта по коэффициенту Коэна.
    \item[$N$] --- число одновременно подключённых участников при нагрузочном тестировании.
    \item[$p_{50}, p_{95}, p_{99}$] --- 50-й, 95-й и 99-й перцентили латентности.
\end{description}
